\section{Complete estimation budgets for a certified financial correction}
\label{sec:integral-estimation-budget}

I retain the established canonical amplitude-estimation theorem
\citep{brassard}. The new financial normalisation can require a probability
tolerance below the earlier implementation domain $10^{-6}$, so I extend
the resource implementation without silently extrapolating its tapes.
The construction in this section supports phase widths through 64. Every
accepted width has a retained circuit and a norm bound; a request beyond
that domain is rejected.

\subsection{An explicitly charged wider inverse Fourier transform}

A controlled phase of distance $d$ has matrix
$\operatorname{diag}(1,1,1,e^{-i\pi/2^d})$. Replacing it by the identity
has operator-norm error
\[
 2|\sin(\pi/2^{d+1})|\leq\pi/2^d.
\]
The inequality follows from $|\sin x|\leq|x|$, which follows by integrating
$\cos x$ and using $|\cos x|\leq1$. In an $m$-wire Fourier circuit there
are $m-d$ controlled phases of distance $d$, for $1\leq d<m$.
I emit the existing certified sequences for $d\leq31$ and replace more
distant phases by the identity. Hence the extra omission allowance is
\begin{equation}
 \Delta_m^{\rm omit}\leq\pi_+
       \sum_{d=32}^{m-1}{m-d\over2^d},\qquad \pi_+=355/113>\pi.
 \label{eq:integral-qft-omission}
\end{equation}
The sum is zero for $m\leq32$. All retained phases keep their scalar phase
corrections. An omitted phase has zero scalar correction; it is compared
directly with the identity as above.

For unitary products, insertion and subtraction of one factor at a time
gives an operator-norm difference bounded by the sum of individual
factor errors, since multiplication by a unitary preserves that norm.
Combining the retained rotation certificates with
\eqref{eq:integral-qft-omission} therefore gives a complete bound
$\Delta_m$. The aggregate scalar correction can be commuted to the front
of the product. It has no effect on this final, uncontrolled transform's
measurement probabilities. It cannot be discarded in a subsequently
controlled use, which is not the construction here.

The implemented widths one through 32 reproduce the previous exact
rotation tapes, gate counts and norm bounds. The extension retains actual
chronological gate sequences and their hashes for every wider profile.
At width 64 the certified norm upper bound is approximately $0.002202$,
which is below $1/64$. Its final-transform cost is 167,223 T gates.
These counts concern the Fourier component, rather than the financial
preparation or a complete pricing run.

\subsection{Confidence and the monetary target}

Let $p$ be the exact probability of a certified finite flag and let
$u\in[0,1]$ be a certified upper prior. Set
$v=u(1-u)$ when $u\leq1/2$ and $v=1/4$ otherwise. A rational upper
square root of $v$ supplies the usual sufficient amplitude-estimation
radius
\begin{equation}
 e_M={2\pi_+\sqrt{v}^{\,+}\over M}+{\pi_+^2\over M^2},
 \qquad M=2^m.
 \label{eq:integral-canonical-radius}
\end{equation}
I choose the first supported $m$ for which
$e_M+10^{-60}\leq e_p$, where $10^{-60}$ is the retained, separately
proved classical output-arithmetic allowance. I require
$10^{-60}<e_p<1/2$. If $u\leq2e_p$, the midpoint of $[0,u]$ already
suffices without shots. This is distinct from the financial whole-correction
or continuous-pricing tests, which can resolve a price before a flag is
constructed at all.

To transfer the compiled Fourier error to its confidence, let $\psi$ and
$\widetilde\psi$ be normalised ideal and compiled joint states, with their
scalar phases aligned and distance at most $\Delta_m$. Write
$c=\langle\psi,\widetilde\psi\rangle$. The difference of their rank-one
projectors has trace zero and its two possible nonzero eigenvalues are
$\pm\sqrt{1-|c|^2}$: its squared trace is $2(1-|c|^2)$ and its rank is at
most two. For any event projector, its expectation difference consequently
has magnitude at most $\sqrt{1-|c|^2}$. Moreover
\[
 1-|c|^2\leq2-2\operatorname{Re}c
       =\|\psi-\widetilde\psi\|^2,
\]
because the difference between the right and left sides is $|1-c|^2$.
Thus an ideal good-event probability lower bound $8/\pi_+^2$ becomes
$8/\pi_+^2-\Delta_m$. This reasoning includes entanglement with the
financial registers; no pure reduced-state assumption is required.

I round $\Delta_m$ upward to the $2^{-80}$ grid and call the resulting
success lower bound $s_m$. It is strictly above one half for every retained
profile. Independent fresh repetitions and an odd median give failure at
most the exact finite sum
\begin{equation}
 \sum_{j=(n+1)/2}^n\binom nj(1-s_m)^j s_m^{n-j}.
 \label{eq:integral-confidence-binomial}
\end{equation}
For success probabilities at least $s_m$, a common-uniform coupling makes
the number of failures stochastically no larger than a binomial law with
parameter $1-s_m$, proving the bound. I take the first odd $n$ whose sum
is at most the prescribed $\alpha$. The phase-width search and confidence
search are finite under the accepted-domain conditions. A reproducible
simulation seed does not establish independent physical repetitions.

For a signed correction $C\mathbb E X$, decoded by $2p-1$, the estimation
radius is $2C e_p$. For a nonnegative correction decoded by $p$, it is
$C e_p$. Here $C$ includes any Gaussian-likelihood normalisation of the
chosen preparation. Coefficient, financial transfer, auxiliary-grid and
arithmetic errors consume their stated monetary allowances before choosing
$e_p$. They are not charged again as a statistical failure event.
For 27 prespecified selected prices I allocate $\alpha=1/2700$ to each;
the union bound gives total failure at most $1/100$. The implementation
selection is deterministic and precedes all estimator shots. This coverage
is for one selected implementation per price, rather than simultaneous
coverage of every tested method and precision.

\subsection{Full coherent operations and reusable-shot accounting}

Let $A$ be the complete financial comparator preparation, including its
original-label Hadamards, auxiliary calculation, numerical evaluation,
flag and private-work erasure. The reflection core consists of original
input labels, comparison labels and the flag; all excluded private work is
zero on its invariant subspace. The controlled Grover circuit is the
chronological sequence
\[
 \operatorname{CZ}_{C,F},\quad A^\dagger,\quad
       \operatorname{controlled}S_0,\quad A,\quad Z_C.
\]
The final $Z_C$ supplies the conditional global minus. Arithmetic gates in
$A$ are not all upgraded to controlled gates. The zero reflection computes
its conjunction with actual Toffolis and returns its private work to zero.
I expand each Toffoli by its literal seven-T sequence. The supplied
financial profile must already consist of coherent Clifford and T atoms;
a remaining Toffoli placeholder or measurement gadget is rejected.

The primitive-count record for $A^\dagger$ swaps T with its adjoint and S
with its adjoint, while leaving self-adjoint gates unchanged. Doubling the
individual T count would misreport its actual gate types even though the
sum of T and its adjoint is unchanged. The complete cost per repetition
is one $A$, $M-1$ controlled Grover calls, phase-input Hadamards, one emitted
inverse Fourier tape, final observations and reset. Serial T-depth caps,
logical-qubit allocations and operation counts are kept separately.

For a reusable shot with $w$ allocated wires, I observe each wire and apply
its classically conditioned X reset. This uses $w$ observations and at
most $w$ reset operations. Phase-register observations are already present
in the estimator's observation metric, so I add only $w-m$ further
observations there. They were not present in its primitive-operation metric;
that metric receives all $2w$ observation/reset operations exactly once.
Both the per-repetition and aggregate records contain these charges.
Each total equals its per-repetition count times the number of repetitions.
Reset has no T cost, but its operations and allocated wires are retained.

The circuit profiles and confidence bounds specify an ideal logical
implementation. They do not establish physical execution time, a quantum
noise threshold or superiority over classical pricing. A financial method
enters this comparison only after its own continuous-target and complete
value-oracle certificates have passed.
